\documentclass[12pt, a4paper]{article}

\usepackage{amsmath}
\usepackage{amssymb}
\usepackage{amsthm}
\usepackage{graphicx}
\usepackage{geometry}
\usepackage{parskip}
\usepackage{microtype}

\geometry{margin=2.5cm}

\title{Converted Notes}
\author{Reuben Stannah}
\date{\today}
\begin{document}
\maketitle
\section{MDM Bullet Points}
\subsection{Slide 11}
- This in a model aiming to emulate the physics of the metronomes. \\
- 4 assumptions; List all of them \\
- First, consider a single metronome on a movable platform \\
- Modeled as a Van Der Pol oscillator. \\
- $\epsilon$ is a perimeter that encodes a non-linear term proportional to the square of the speed of the platform. This is used to decay the oscillations with time. \\


\subsection{Slide 12}
- This in extended slightly into \textit{the system} of 2 equations of motion. \\
- Note that Assumption 4 $\Rightarrow$ the equation of motion of the centre of mass of the system may be found using Newton's $2^{nd}$ Law. \\
- Full deviation is the technical note but the equation of motion is given by:

\subsection{Slide 13}

- The case \( n=2 \) oscillator model may be extended to an arbitrary amount \\
- The state vector may be defined as g, such that encoding all angles $\theta_1 \dots \theta_n$ as well as x. \\
- M, G, \& C are matrices where elements are functions of q and its derivatives. \\
- This allows whole system to be written using the equation \( S = f(\mathbf{q}, \vdot{\mathbf{q}}) \). \\
- Hence numerical methods, we used RK4 for its simplicity and flexibility, may be used to investigate and plot the effects of different control strategies. \\

\subsection{Slide 14}

- These are M, C, G and T; full derivations in handwritten note. \\

\end{document}